\documentclass[11pt,a4paper]{article}
\usepackage{fontspec}
\usepackage{polyglossia}
\setmainlanguage{russian}
\setotherlanguage{english}
\setmainfont{Times New Roman}
\usepackage{booktabs,geometry,array,longtable}
\geometry{margin=2.2cm}
\usepackage{hyperref}
\hypersetup{hidelinks}

\title{Таргетированные состязательные атаки\\на детекторы синтетических изображений}
\author{Отчёт по результатам кампании}
\date{4 сентября 2026 г.}

\begin{document}
\maketitle

\section{Задача}

Проверялось, можно ли незаметным изменением пикселей заставить детектор
дипфейков ошибиться, и переносится ли такая атака на другой детектор, который
она не видела.

Второй вопрос главный. Атака, работающая только против той модели, по которой
она строилась, говорит о слабости конкретной сети. Атака, переносящаяся между
моделями разной архитектуры, указывает на общую уязвимость.

\section{Постановка}

Два детектора: \textbf{ViT-B/16} работает с пространственным RGB-изображением,
\textbf{DenseNet121-DCT} — с частотным представлением.

Каждый детектор выступал в двух ролях. \emph{Источник} — модель, по градиенту
которой строится возмущение. \emph{Цель} — модель, на которой проверяется
результат. Совпадение ролей даёт белый ящик, различие — перенос.

Данные: 100 поддельных изображений из \texttt{TEST/TEST\_FAKE} набора CelebA,
замороженный манифест с контролем хеша. Направление атаки — fake${\to}$real,
то есть подделку заставляют выглядеть настоящей.

Ограничение возмущения: $L_\infty \le 8/255$, проверяется по перечитанному с
диска PNG. Сид 0 во всех прогонах.

Метрики: \textbf{ASR} — доля успеха, где знаменатель это изображения,
классифицированные детектором верно \emph{до} атаки; \textbf{SSIM} и
\textbf{LPIPS} — сохранность изображения.

\section{Атаки}

Шесть методов, от простого к сложному.

\begin{longtable}{@{}lp{11.2cm}@{}}
\toprule
\textbf{Метод} & \textbf{Суть} \\
\midrule
\endhead
FGSM & Один шаг по знаку градиента на весь бюджет. \\
PGD & Десять малых шагов со случайным стартом и проекцией после каждого. \\
MI-DI-FGSM & PGD плюс momentum и случайное изменение размера входа. \\
Ensemble MI + EoT & Momentum по обоим детекторам, усреднение по искажениям. \\
SSA / S$^2$I-FGSM & Усреднение градиентов по 20 случайно зашумлённым спектрам. \\
MIG-COW & Интегрированные градиенты с консенсусным взвешиванием. \\
\bottomrule
\end{longtable}

\section{Результаты}

Двенадцать прогонов, по 100 изображений каждый. Все прошли детерминированную
верификацию, нарушений бюджета ноль.

\begin{longtable}{@{}llrrrrr@{}}
\toprule
\textbf{Метод} & \textbf{Источник} & \textbf{Бел. ящик} & \textbf{Перенос} &
\textbf{SSIM} & \textbf{LPIPS} & \textbf{с} \\
\midrule
\endhead
FGSM              & ViT & 0.010 & 0.010 & 0.864 & 0.303 & 178 \\
FGSM              & DCT & 0.160 & 0.000 & 0.921 & 0.099 & 177 \\
PGD               & ViT & 0.780 & 0.040 & 0.848 & 0.211 & 182 \\
PGD               & DCT & 0.250 & 0.000 & 0.815 & 0.197 & 154 \\
MI-DI-FGSM        & ViT & 0.130 & 0.010 & 0.907 & 0.246 & 167 \\
MI-DI-FGSM        & DCT & 0.140 & 0.000 & 0.935 & 0.100 & 205 \\
Ensemble MI + EoT & ViT & 0.280 & 0.020 & 0.914 & 0.234 & 252 \\
Ensemble MI + EoT & DCT & 0.240 & 0.000 & 0.948 & 0.077 & 274 \\
MIG-COW           & ViT & 0.720 & 0.010 & 0.907 & 0.241 & 581 \\
MIG-COW           & DCT & 0.290 & 0.000 & 0.952 & 0.072 & 432 \\
SSA               & ViT & \textbf{0.870} & 0.040 & 0.890 & 0.273 & 662 \\
SSA               & DCT & \textbf{0.320} & 0.000 & 0.765 & 0.375 & 985 \\
\bottomrule
\end{longtable}

\subsection*{Матрица переноса}

Строки — источник градиента, столбцы — где измеряли. Приведён лучший результат
по каждой клетке.

\begin{center}
\begin{tabular}{@{}lrr@{}}
\toprule
& \textbf{ViT} & \textbf{DCT} \\
\midrule
\textbf{ViT} (источник) & 0.870 & 0.040 \\
\textbf{DCT} (источник) & 0.000 & 0.320 \\
\bottomrule
\end{tabular}
\end{center}

\section{Выводы}

\textbf{Перенос отсутствует.} Это главный результат. Лучшее значение вне
диагонали — 0.040, то есть четыре изображения из ста, и достигается оно двумя
разными методами. В направлении DCT${\to}$ViT перенос ровно нулевой у всех
шести атак без исключения.

Ни один метод не изменил картину, включая те, что специально проектировались
ради переносимости: momentum и диверсификация входа в MI-DI, усреднение по
спектрам в SSA, ансамблирование и консенсус в MIG-COW. Пространственный и
частотный детекторы, судя по этим данным, не разделяют эксплуатируемых
уязвимостей.

\textbf{SSA сильнейшая в белом ящике} — 0.870 против ViT и 0.320 против DCT,
лучший результат в обоих направлениях. Работа в частотной области помогает
атаковать, но не переносить.

\textbf{Атака через DCT слабее.} При одинаковом бюджете белый ящик по частотному
детектору держится в диапазоне 0.14–0.32 против 0.01–0.87 по пространственному.
Вероятная причина в том, что градиент проходит через перевод в градации серого,
масштабирование, обрезку до $128{\times}128$ и DCT, теряя большую часть пути к
пикселям.

\textbf{FGSM практически не работает} — 0.010 в белом ящике при худшем качестве
изображения. Один грубый шаг на весь бюджет промахивается там, где десять малых
с проекцией попадают.

\section{Оговорки}

\textbf{MI-DI-FGSM ведёт себя аномально.} Метод есть PGD с добавленными
momentum и диверсификацией, то есть должен быть не слабее. Фактически 0.130
против 0.780 у PGD, в шесть раз хуже. Это указывает на дефект реализации,
вероятнее всего в нормировке градиента перед накоплением момента, а не на
свойство метода. Строку следует считать недостоверной до перепроверки.

\textbf{Ансамблевые методы вырождены.} При двух детекторах схема
leave-one-detector-out оставляет один источник. Ensemble MI + EoT и MIG-COW
запускались, но их результаты нельзя читать как ансамблевые.

\textbf{Конфаунд по формату файлов.} В наборе данных все настоящие изображения
хранятся в JPG, все поддельные — в PNG. Метку можно предсказать по расширению
файла, не анализируя пиксели. Чистая точность обоих детекторов равна 1.000, и
она может отражать различие в сжатии, а не способность распознавать подделку.
Это ограничение самого набора данных.

\textbf{Отложенной выборки нет.} Все методы сравнивались на одних и тех же 100
изображениях, на них же приведены результаты. Параметры атак задавались заранее
в спецификациях и не подбирались по ответам детекторов, что ограничивает
подгонку, но строгое утверждение об обобщении требует непересекающегося набора.

\textbf{Матрица неполна по детекторам.} Веса детекторов NPR и AIDE доступны, но
адаптеры не написаны. С двумя моделями перенос измеряется единственной парой.

\section{Что дальше}

Первое — перепроверить MI-DI-FGSM: если после исправления он окажется на уровне
PGD, картина по переносу не изменится, но станет непротиворечивой.

Второе — подключить NPR. Третий детектор делает осмысленными и ансамблевые
методы, и leave-one-detector-out, и превращает единственную пару в матрицу.

Третье — устранить конфаунд по формату, приведя все изображения к единому
контейнеру. Без этого любое утверждение о точности детекторов остаётся под
вопросом.

\vfill
\footnotesize
Все числа получены из первичных файлов \texttt{summary.json}, каждый из которых
сопровождается отчётом верификатора с \texttt{outcome=passed}. Воспроизводимость
подтверждена: два идентичных прогона дали побайтово совпадающие изображения.

\end{document}
